% !TEX root = ../Main.tex
\chapter{发球与接球}

\section{基本站位与角色}

\state{双方的站位}{
    竞赛双方站在球网两边，其中一名（或一组）运动员为\textbf{发球手}，另一位（或一组）则是\textbf{接球手}。\textbf{一局比赛结束之后交换发球权}。

    \begin{itemize}
        \item 发球手必须站在\textbf{底线后面}，以及\textbf{中界点和边线}之间的任意一点；
        \item 接球手可选择站在\textbf{球网后面}的任何一点，一般站在发球手的\textbf{斜对面}；
        \item 发球时，选手应该用手将球\textbf{抛向空中}任何方向，然后在球\textbf{落地之前}用球拍击球。
    \end{itemize}
}

\section{发球动作要求}

\state{发球细则}{
    \begin{itemize}
        \item 如果球手对抛球动作不满意，可以让球\textbf{自然落地}，然后重新开始；
        \item 如果球手做出用球拍击球的动作但\textbf{并未击打到球}，则被视为\textbf{发球失误}；
        \item 发球手在发球的整个过程中，\textbf{两脚不应偏离最初的位置}；发球手的脚\textbf{可以离开地面}，但是\textbf{不得走动或奔跑}。
    \end{itemize}
}

\section{发球是否有效}

\state{有效发球的条件}{
    在一次正确的发球中，发出的球应该：
    \begin{itemize}
        \item \textbf{越过球网}（但不触及球网）；
        \item \textbf{落在球网另一侧斜对角发球区中的任意一点}。
    \end{itemize}
}

\state{特殊情形}{
    \begin{itemize}
        \item \textbf{触网（Net）}——球\textbf{触网但依然落在斜对角的发球区}：\textbf{重新发球}（\textbf{不记为发球失误}）。
        \item \textbf{发球失误}——球\textbf{落网 / 没过网前就落地 / 落地前触碰除球网外的任何物体 / 没有落在斜对角发球区}。
        \item \textbf{双误（Double Fault）}——第一次发球失误后发球手可以再发一次，如果\textbf{依然失误}，让对手得一分。
    \end{itemize}
}

\section{接球规则}

\state{接球要求}{
    \begin{itemize}
        \item 球手在球\textbf{落地回弹之前}就击球，依然有效，此为"\textbf{截击}"（Volley）；
        \item 以下行为将导致\textbf{失分}：
            \begin{itemize}
                \item 触碰球网；
                \item 在球\textbf{穿越球网前}击球；
                \item 用\textbf{球拍之外}的任何物体击球；
                \item \textbf{故意击球两次}等。
            \end{itemize}
        \item 在\textbf{双打}比赛中，在发球与首次回击后，\textbf{搭档中的任何一位在任意回合中都可击球}。
    \end{itemize}
}

\rmkb{
    \textbf{易考点}：
    \begin{itemize}
        \item 发球要\textbf{在球落地之前}击打；
        \item \textbf{触网但进发球区}$\to$ \textbf{重发}，不记失误；
        \item \textbf{双误}：\textbf{连续两次发球失误} $\to$ 对手得分；
        \item \textbf{截击}：\textbf{球落地前击球}；
        \item 发球时\textbf{脚可离地但不可走动}。
    \end{itemize}
}
